\documentclass[french,12pt,a4paper]{article}
\usepackage[T1]{fontenc}
\usepackage[width=2.5cm, height=2.5cm, left=2.5cm, right=2.5cm, top=2.5cm, bottom=2.5cm]{geometry}
\usepackage[utf8]{inputenc}
\usepackage{amsmath, amssymb}
\usepackage{babel}
\usepackage{graphicx}
\usepackage{tcolorbox}
\usepackage{xcolor}
\begin{document}
   	\begin{tabular}{p{3cm}p{11cm}p{3cm}}
   		&\begin{center}
  	\includegraphics[width=7cm,height=0.4cm]{../../../Desktop/drapeau (1)}
   		\end{center}&\\
   		\vspace{0.3cm}
   		\includegraphics[width=3cm,height=4cm]{../../../Desktop/ois}	& 
   		\begin{center}
   			\textbf{\large République du Bénin}\\
   			\hspace{-0.5cm}\textbf{\large Ministère de l'Enseignement }%\\[0.1cm]
   			\textbf{\large Supérieur et de la   Recherche  }%\\[0.1cm]
   			\textbf{\large Scientifique (MESRS)}\\ \vspace{0.22cm}
   			
   			\textbf{ \large	Université Nationale des Sciences, Technologies, Ingénierie et Mathématiques (UNSTIM) d'Abomey} \\ [0.1cm]
   			
   			\textbf{ÉCOLE NORMALE SUPÉRIEURE DE NATITINGOU}
   		\end{center}
   		& \vspace{0.5cm} \includegraphics[width=2cm,height=3cm]{../../../Desktop/ENS}
   	\end{tabular}	
	 	\noindent\textcolor{blue}{\rule{17cm}{0.05cm}}
	\section*{I. MATHEMATIQUES (7 pts)}
	 	\noindent\textcolor{blue}{\rule{17cm}{0.05cm}}
	Soit $f$ une fonction continue périodique. On considère ses coefficients :
	\[	c_n(f)=\frac{1}{2\pi}\int_{0}^{2\pi} f(t)e^{-int}\,dt	\]
	\begin{enumerate}
		\item Prouver que :
		\[
		\sum_{n=-\infty}^{+\infty} |c_n(f)|^2
		= \frac{1}{2\pi}\int_{0}^{2\pi} |f(t)|^2 dt
		\]
	\end{enumerate}	
	\section*{II. PHYSIQUE (7 pts)}
	 	\noindent\textcolor{yellow}{\rule{17cm}{0.05cm}}
	Démontrer l’invariant de la pseudo-norme du quadrivecteur impulsion :
	\[\left(\frac{E}{c}\right)^2 - \vec{p}^{\,2} = m^2 c^2\]
	\section*{III. CHIMIE (6 pts)}
	 	\noindent\textcolor{green}{\rule{17cm}{0.05cm}}
	Calculer l’enthalpie libre de réaction $\Delta_r G$ en fonction du quotient de réaction $Q$ :
	\[\Delta_r G = \Delta_r G^\circ + RT \ln(Q)	\]
 
	\colorbox{yellow}{\parbox{\linewidth}{En Afrique subsaharienne ,la transformation digitale en santé se dépoloie dans un contexte singulier . Les systèmes de santé y sont caractérisés par une insuffisance chronique de ressources financières ,humaines et matérielles (Nabyongo-Orem \&  Tashobya ,2016).Les infrastructures numériques,bien que croissantes avec l'essor de la téléphonie mobile et d'Internet,restent inégales et parfois précaires . Les coupures d' électricité fréquentes , l'instabilité des réseaux et la faible capacité technique du personnel constituent des contraintes majeures (Dovlo,2018). Cependant, ces fragilités coexistent avec un réel potentiel: }}
	
  
\end{document}

